\documentclass{situationist-preprint}

\title{Microphysical Consciousness and the Combination Problem}
\msid{MND-01-d}
\msfamily{MND-01}
\msdiscipline{Philosophy of Mind}
\mstype{Consequence}
\mskeywords{Russellian monism; protophenomenal properties; combination problem; boundary problem; palette problem}

\begin{document}
\maketitle

\begin{abstract}
If consciousness is the intrinsic nature of physical structure, are the fundamental entities of physics themselves conscious? We derive from the conclusion of MND-01, together with the definition of the physical defended in MND-02 and the result of MND-18, that they are not: their intrinsic natures are protophenomenal. It follows that the combination problem in its strongest form, the problem of summing subjects, does not arise for the view. Two problems remain, the boundary problem and the palette problem. We give a preliminary treatment of each and argue that both are problems of stating the principles of a grounding relation, not of showing that the relation can exist.
\end{abstract}

\section{No microphysical consciousness}\label{s:micro}

MND-01 holds that the hard problem is solved because consciousness is the intrinsic nature of physical structure; given materialism, Russellian monism is true (MND-01-a, Corollary~4.3). The formula leaves open where in physical structure consciousness lies. Panpsychist versions place it at the bottom: every fundamental entity is a subject, and our experience is grounded in the experiences of such subjects \citep{Strawson2006,Goff2017}. Protophenomenal versions deny this: fundamental natures are not experiences and do not make their bearers subjects, though in suitable combination they ground experiences (MND-18-a, Definition~2.3; \citealp{Chalmers2013}).

\begin{proposition}\label{p:nomicro}
On the Russellian monism of MND-01, fundamental physical entities are not conscious, and their intrinsic natures are protophenomenal.
\end{proposition}

\begin{proof}
Two independent routes lead to the same result. First, MND-01 takes materialism in the sense of MND-02, on which a property is physical only if it is a fundamental property of a foreseeable extension of current physics, or is grounded in such properties; and no fundamental property of a foreseeable extension is mental (MND-02-a, Definitions~2.2(c) and 2.3). If the natures that MND-01 posits are to be physical, as Corollary~4.3 of MND-01-a requires, they are not mental, and so are not experiences. Second, MND-18 shows that if consciousness is the intrinsic nature of physical structure, the intrinsic natures of fundamental entities are not experiences of those entities (MND-18-a, Corollary~3.3), since otherwise our experience would be grounded in the experiences of other subjects, which its Theorem~3.2 excludes. Natures that are not experiences, but ground experiences in combination, are protophenomenal.
\end{proof}

The title question therefore has a negative answer. An electron has an intrinsic nature, and that nature is of the same general kind as what is presented to us intrinsically in experience, in that structures composed of such natures are conscious; but the electron is not a subject, and there is nothing it is like to be one.

\begin{corollary}\label{c:nosum}
The combination problem in the form of subject-summing does not arise for the Russellian monism of MND-01.
\end{corollary}

\begin{proof}
The subject-summing problem asks how the experiences of many subjects can ground the experience of a further subject. By Proposition~\ref{p:nomicro} there are no micro-subjects, so the grounding of our experience is not grounding in the experiences of others; by Proposition~4.1 of MND-18-a, the principles that make subject-summing impossible are consistent with grounding of this kind.
\end{proof}

\section{The residual problems}\label{s:residual}

What remains is the task of saying which arrangements of protophenomenal natures ground experiences, and which experiences they ground \citep{Chalmers2017}. Two parts of the task have names.

\emph{The boundary problem.} Protophenomenal natures are everywhere in the brain and outside it. Why does their arrangement ground exactly one subject, bounded roughly by the boundaries of a cortical system, rather than many overlapping subjects or none? Our preliminary answer is that the grounding principle is maximal: an arrangement grounds a subject only if it realizes the relevant integrated causal organization and is not a proper part of an arrangement that realizes it. Overlapping candidates are then excluded by construction, and the boundaries of a subject are the boundaries of a maximal integrated arrangement. Which organization is relevant is an empirical question for the science of consciousness; the principle fixes only the form of the answer.

\emph{The palette problem.} Our experience presents a great variety of qualities, while physics posits a small number of fundamental properties. How can a few natures ground many qualities? Here the structural account of qualia in MND-10 does part of the work. The structure of a quality space is fully describable relationally (MND-10-a, Theorem~5.4), and the variety of relational structures that arrangements of a few natures can realize is unbounded, as the variety of chemical compounds realized by a few kinds of particle shows. What remains is the dependence of the intrinsic character of each macro-quality on the natures in its ground. Our preliminary proposal is that this dependence is fixed by grounding principles that pair the intrinsic character of a quality with the natures and arrangement of its realizer, just as the intrinsic character of the realizer is paired with its structure in MND-01.

\begin{proposition}\label{p:articulation}
Neither the boundary problem nor the palette problem, if unsolved, yields an inconsistency in protophenomenal Russellian monism. Each is a problem of stating the principles of the grounding relation, not of showing that the relation can exist.
\end{proposition}

\begin{proof}
An inconsistency would require that some principle about subjects or experiences be violated by any grounding of experience in protophenomenal natures. MND-18 identifies the only principles known to generate such a conflict, Inclusion, Ownership and Exclusivity, and shows them consistent with protophenomenal grounding (MND-18-a, Proposition~4.1). The boundary problem asks which arrangements ground subjects; the palette problem asks which qualities they ground. Both presuppose that some arrangements ground some qualities and ask which, as the special composition question presupposes that some pluralities compose something and asks which \citep{vanInwagen1990}. An unanswered question of this form is a gap in the theory, not a contradiction in it.
\end{proof}

\begin{objection}
Protophenomenal properties are characterized only negatively, as not experiences but able to ground them. A theory that posits them explains nothing about our experience until it says what they are.
\end{objection}

\begin{reply}
They are characterized positively in one respect: they are the intrinsic natures that physics leaves unspecified, and in experience we are acquainted with what some arrangements of them are like. MND-01 locates the explanatory gap in the difference between this intrinsic mode of presentation and the structural mode of physics; it does not require an independent conception of the natures one by one. The demand for one is a demand for the complete grounding principles, which is the task Volume~IV of MND-01-a takes up.
\end{reply}

\section{Conclusion}

On the Russellian monism of MND-01, consciousness is the intrinsic nature of some physical structures but not of the fundamental entities of physics, whose natures are protophenomenal. The combination problem in its fatal form does not arise; what remains are the boundary and palette problems, which ask for the principles of a grounding relation and are developed in Volume~IV of MND-01-a.

\begin{thebibliography}{9}\small
\bibitem[Chalmers(2013)]{Chalmers2013} Chalmers, D.~J. (2013). Panpsychism and panprotopsychism. \emph{The Amherst Lecture in Philosophy}, 8.
\bibitem[Chalmers(2017)]{Chalmers2017} Chalmers, D.~J. (2017). The combination problem for panpsychism. In G.~Br\"untrup and L.~Jaskolla (eds.), \emph{Panpsychism: Contemporary Perspectives}. Oxford University Press.
\bibitem[Goff(2017)]{Goff2017} Goff, P. (2017). \emph{Consciousness and Fundamental Reality}. Oxford University Press.
\bibitem[Strawson(2006)]{Strawson2006} Strawson, G. (2006). Realistic monism: why physicalism entails panpsychism. \emph{Journal of Consciousness Studies}, 13(10--11), 3--31.
\bibitem[van Inwagen(1990)]{vanInwagen1990} van Inwagen, P. (1990). \emph{Material Beings}. Cornell University Press.
\end{thebibliography}

\end{document}
